\chapter{Design}
\label{cap:design}

\section{Class Diagram}
\label{sec:class_diagram}

\subsection{VOPC (Analysis)}
\label{subsec:vopc}

\begin{figure}[H]
    \centering
    \includegraphics[width=\textwidth]{figures/vopc_uc04.png}
    \caption{VOPC di UC-04: Boundary (View), controller grafico, Facade, controller applicativi ed Entity.}
    \label{fig:vopc}
\end{figure}

\subsection{Design Patterns}
\label{subsec:design_patterns}
Il design-level introduce le interfacce, i tipi Java, le eccezioni e i pattern
GoF per gestire estensibilità e persistenza dinamica (NFR-01). Per ogni pattern
sono mostrati in ordine lo schema e la breve descrizione dell'applicazione.

\subsubsection*{Abstract Factory Pattern -- Persistence Layer}
\begin{figure}[H]
    \centering
    \includegraphics[width=0.85\textwidth]{figures/dl_factory_hr.png}
    \caption{Abstract Factory della persistenza: interfacce DAO, concrete DAO e factory.}
    \label{fig:dl_factory}
\end{figure}
Il livello di persistenza è gestito tramite il pattern \textbf{Abstract Factory}
al fine di garantire flessibilità di configurazione dell'ambiente di esecuzione.
L'interfaccia \texttt{DAOFactory} definisce il contratto per la creazione dei vari
Data Access Object, come \texttt{getOrdineDAO()} e \texttt{getProdottoDAO()}.
In base alla proprietà \texttt{persistence\_type}, le factory concrete
(\texttt{JDBCDAOFactory}, \texttt{FSDAOFactory}, \texttt{DemoDAOFactory})
istanziano l'opportuna implementazione dello storage. Ciò consente al sistema
di passare senza problemi tra una persistenza basata su database, file-system
o memoria in-memory senza richiedere alcuna modifica al codice client.

\subsubsection*{Observer Pattern -- Notification System}
\begin{figure}[H]
    \centering
    \includegraphics[width=0.55\textwidth]{figures/dl_observer_hr.png}
    \caption{Observer della conferma dell'ordine: publisher singleton e DTO.}
    \label{fig:dl_observer}
\end{figure}
Il pattern \textbf{Observer} è stato implementato con lo scopo di disaccoppiare
l'elaborazione dell'ordine dalla logica di gestione delle notifiche. La classe
\texttt{OrdineEventPublisher} funge da \emph{Subject} implementato tramite il
pattern \textbf{Singleton}, gestendo una lista (\texttt{List}) di osservatori
di tipo \texttt{OrdineEventListener}. Al momento del corretto invio e della
conferma dell'ordine, il publisher attiva il metodo \texttt{onOrdineConfermato()}
su tutti i listener registrati (compratore e venditore), passando il DTO
\texttt{OrdineEvent} in sola lettura, mai l'entità di dominio. Questa
architettura garantisce un accoppiamento debole tra la logica di business
principale e i relativi effetti collaterali, rispettando rigorosamente
l'\emph{Open/Closed Principle}.

La notifica \`e emessa dal sottosistema applicativo (il control che conferma
l'ordine) tramite il publisher, e non dall'entity \texttt{Ordine}: cos\`{}\ l'entit\`a
di dominio resta \emph{svincolata} dal layer di presentazione e dal meccanismo
di notifica (separazione tra interfaccia grafica e struttura dati indicata in
teoria). La variante \`e quella \emph{publish-subscribe}/event-driven
dell'Observer: l'accopiamento tra \texttt{OrdineEventPublisher} (Subject) e gli
observer \`e astratto (solo \texttt{OrdineEventListener}), la notifica \`e un
broadcast a tutti i listener registrati e il dato scambiato \`e il DTO
\texttt{OrdineEvent} in sola lettura, mai l'entit\`a.

\subsubsection*{Singleton Pattern -- Sessions and Mode}
\begin{figure}[H]
    \centering
    \includegraphics[width=0.55\textwidth]{figures/dl_singleton_hr.png}
    \caption{Singleton della sessione e della modalità.}
    \label{fig:dl_singleton}
\end{figure}
\texttt{SessionManager} (sessione utente) e \texttt{ApplicationModeManager}
(modalità applicativa) adottano un singleton classico con \emph{lazy
initialization} (\texttt{if (instance == null) instance = new X()}) e
sincronizzazione su \texttt{getInstance()}, in linea coi progetti di
riferimento: nessun costo percepibile in un client desktop, nessuna
soppressione di warning per SonarCloud.

\subsubsection*{Facade Pattern -- Checkout UC-04}
\begin{figure}[H]
    \centering
    \includegraphics[width=0.75\textwidth]{figures/dl_facade_hr.png}
    \caption{Facade del checkout di UC-04.}
    \label{fig:dl_facade}
\end{figure}
\texttt{OrdinaProdottoFacade} espone al controller grafico un'unica interfaccia
per il checkout (\texttt{getProdottiDisponibili},
\texttt{avviaCheckout}, \texttt{applicaBuono}, \texttt{processaPagamento})
e nasconde il sottosistema di business, formato dal controller base e dai
controller delle estensioni (buono 4a e pagamento 6/6a). Il client (boundary
o controller grafico) non conosce l'ordine di chiamata ai metodi dei
controller applicativi.

\section{Activity Diagram}
\label{sec:comportamentale}

\begin{figure}[H]
    \centering
    \includegraphics[height=0.8\textheight]{figures/activity_uc04.png}
    \caption{Activity diagram di UC-04: verifica disponibilità, applicazione
    del buono, autorizzazione del pagamento e notifica sequenziale a
    compratore e venditore.}
    \label{fig:activity}
\end{figure}

\section{Sequence Diagram}
\label{sec:sequence}

\begin{figure}[H]
    \centering
    \includegraphics[width=\textwidth]{figures/seq_uc04_clean.png}
    \caption{Sequence diagram di UC-04.}
    \label{fig:sequence}
\end{figure}

\section{State Diagram}
\label{sec:state}

\begin{figure}[H]
    \centering
    \includegraphics[width=\textwidth]{figures/state_uc04.png}
    \caption{State diagram di UC-04 Ordina un Prodotto: navigazione e flussi alternativi \texttt{3a}, \texttt{4a} e \texttt{6a}.}
    \label{fig:state}
\end{figure}

